\section{Introduzione al Test}
Il Test valuta la capacità di scegliere un dispositivo e una procedura di
compilazione per un circuito quantistico mai usato per scegliere le impostazioni
del sistema. Confrontiamo cinque soluzioni: LLM + RAG, lo stesso LLM senza
esempi (LLM no RAG), MQT Predictor e una scelta casuale, indicata come Random. Si aggiunge LLM + Random RAG: lo stesso LLM riceve
cinque esempi train estratti a caso, anziché recuperati per somiglianza.
L'obiettivo è capire se gli esempi di compilazioni precedenti aiutano il modello
linguistico a prendere decisioni migliori e quale costo comporta questo aiuto.

La valutazione riguarda 90 circuiti riservati al Test. Gli altri 510 circuiti
del corpus sono stati destinati al train (422) e alla validation (88).
Il train fornisce gli esempi e i dati di apprendimento; la validation ha guidato
la scelta della configurazione LLM. Le impostazioni dei cinque sistemi originali erano già fissate.
La variante con esempi casuali è stata aggiunta dopo la lettura del Test:
il confronto che la include è esplorativo e usa un solo seme di recupero.
Per ciascun circuito e sistema è prevista una sola compilazione.

Osserviamo tre aspetti complementari. Il primo è la \textbf{qualità}, espressa
dallo score \emph{expected fidelity}: un valore fra 0 e 1 che stima la fedeltà
del circuito compilato, dove un valore maggiore è migliore. Il secondo è la
\textbf{riuscita}: quanti circuiti arrivano a una compilazione valida e quanti
terminano con un errore o oltre il tempo disponibile. Il terzo è il
\textbf{costo di esecuzione}: tempo totale, tempo del compilatore e, per gli
LLM, tempo di risposta, token utilizzati e richieste di correzione.
Queste misure vanno lette insieme: uno score elevato sui soli casi riusciti
non descrive da solo l'utilità di un sistema.

Prima del Test, l'aspettativa di lavoro era che LLM + RAG ottenesse score
migliori di LLM no RAG e Random, grazie agli esempi pertinenti. Ci si aspettava
inoltre che MQT Predictor, essendo progettato e addestrato per questo compito,
offrisse il risultato complessivamente migliore, anche nei tempi.
Non era stato anticipato un costo temporale di RAG vicino al doppio di quello
del modello senza esempi. Queste aspettative, qui esplicitate dall'autore
nella revisione del report, sono una chiave di lettura dei risultati;
non costituiscono ipotesi statistiche preregistrate.

La qualità è stimata con proprietà sintetiche dei dispositivi, usando la stessa
metrica di MQT Predictor 2.4.0 per tutti i metodi. Non si tratta di esecuzioni
su un computer quantistico reale. Nelle pagine seguenti descriviamo prima
le impostazioni dei cinque sistemi, poi i risultati in tabelle e grafici,
e infine le conclusioni. L'appendice conserva i valori di ogni circuito.

